\begin{papercard}{nissenbaum2009}
\cf{Stages} \stagetag{S} (definition of privacy as appropriate information flow), \stagetag{H} (what makes a flow harmful: breach of context-relative norms), \stagetag{C} (parameters of a norm as coding dimensions). \emph{Privacy in Context: Technology, Policy, and the Integrity of Social Life} (Stanford University Press, 2010). \textbf{Only the table of contents and the introduction are in the corpus} (22 PDF pages, of which pages 6--22 are the introduction). The file originally held front matter only; the registry entry now points to an OCR copy of a scanned Yale PDF of the table of contents and introduction. Quotes below were checked against the OCR text, which contains recognition errors, so exact wording must be verified in the printed book before it is quoted in the thesis. Chapters 7--9 (the framework itself) are not available here.

\cf{Central claim} The author rejects the view that protecting privacy means restricting access or giving people control: ``What people care most about is not simply restricting the flow of information but ensuring that it flows appropriately'' \pc{nissenbaum2009}{7}{What people care most about is not simply restricting the flow of information but ensuring that it flows appropriately}. The account of appropriate flow is the framework of \emph{contextual integrity}. The introduction says it should explain when people perceive technologies as threats and also gives an approach to evaluating systems.

\cf{Context-relative informational norms} The framework says ``finely calibrated systems of social norms, or rules, govern the flow of personal information in distinct social contexts'' such as education, health care and politics \pc{nissenbaum2009}{8}{finely calibrated systems of social norms, or rules, govern the flow of personal information in distinct social contexts}. These norms ``define and sustain essential activities and key relationships and interests, protect people and groups against harm, and balance the distribution of power'' \pc{nissenbaum2009}{8}{These norms, which I call context-relative informational norms, define and sustain essential activities and key relationships and interests, protect people and groups against harm, and balance the distribution of power}. Information technologies alarm us when they flout these norms, that is, when they violate contextual integrity \pc{nissenbaum2009}{8}{Information technologies alarm us when they flout these informational norms}. Not every change in information flow is a problem: some systems serve context-based values and should be accepted, even celebrated.

\cf{Terminology} The book shifts from ``privacy'' to \emph{personal information flows}, which avoids disputes about the concept and keeps precision \pc{nissenbaum2009}{9}{it frames its claims in terms of personal information flows}. ``Personal information'' means information about an identifiable person, following the EU Directive's definition of personal data \pc{nissenbaum2009}{9}{I use it to mean information about an identifiable person}. Technology is treated as \emph{socio-technical}: the object of study is the device ``construed in terms of key social interdependencies'' \pc{nissenbaum2009}{10}{Rather, the object of study is the device construed in terms of key social interdependencies}, here illustrated with the telephone, RFID and vehicle communications.

\cf{Examples of accepted and resisted flows} In healthcare, monitors and recording devices ``complement attentive, responsive caregivers'' and are accepted as part of high-quality care, while video surveillance in parks, store loyalty programmes, wiretapping and monitoring of online transactions meet suspicion \pc{nissenbaum2009}{11}{In healthcare environments such as hospitals and nursing homes, for example, a plethora of devices such as blood-pressure monitors}. A satisfactory account must explain such contrasts.

\cf{Why not interests or abstract rights} Settling conflicts by competition among interest holders favours the powerful \pc{nissenbaum2009}{13}{The trouble with settling conflicts through brute clashes among interest holders is the advantage it gives to those possessing advantages of power}; the examples are Lotus Marketplace: Households (withdrawn 1991 after public outcry, with no admission of wrongdoing) and Total Information Awareness, which resurfaced in other forms (ChoicePoint, fusion centres). Appeals to universal rights (self-ownership, autonomy, liberty) sit ``in the stratospheres of abstraction'' and clash with each other. Contextual integrity is placed in a middle, social realm: ``Between the ground and the heavens \ldots\ is the realm of the social'' \pc{nissenbaum2009}{15}{Between the ground and the heavens, according to the picture I am}.

\cf{Book outline} Part I surveys technology by three capacities: tracking and monitoring, aggregation and analysis, dissemination and publication \pc{nissenbaum2009}{16}{(2) aggregation and analysis, and (3) dissemination and publication}. Part II criticises predominant approaches (value of privacy, the private--public dichotomy, privacy in public). Part III (table of contents: ``Contexts, Informational Norms, Actors, Attributes, and Transmission Principles'', ``Breaking Rules for Good'', ``Privacy Rights in Context: Applying the Framework'') \pc{nissenbaum2009}{5}{Contexts, Informational Norms, Actors, Attributes, and Transmission Principles} gives the framework. Informational norms govern the flow of information about a subject from one party to another, ``taking account of the capacities (or roles) in which the parties act, the types of information, and the principles under which this information is transmitted'' \pc{nissenbaum2009}{19}{taking account of the capacities (or roles) in which the parties act, the types of information, and the principles under which this information is transmitted among the parties}. Contextual integrity works as a metric: ``preserved when informational norms within a context are respected and violated when they are contravened'' \pc{nissenbaum2009}{19}{preserved when informational norms within a context are respected and violated when they are contravened}. The framework yields a \emph{CI Decision Heuristic} and an \emph{Augmented CI Decision Heuristic} \pc{nissenbaum2009}{19}{which it calls the CI Decision Heuristic and the}.

\cf{Presumption and augmentation} The augmented analysis starts from a presumption in favour of entrenched norms but looks to a context's internal purposes, ends and values as benchmarks; new technologies ``deserve to be embraced when they help achieve important social and context-based ends more effectively'' \pc{nissenbaum2009}{20}{new technologies deserve to be embraced when they help achieve important social and context-based ends more effectively}. Because CI demands appropriate flow, not merely control or secrecy, it is said to explain behaviour that sceptics call paradoxical \pc{nissenbaum2009}{20}{demands appropriate flow and not merely control and secrecy, it predicts the behaviors skeptics cite as paradoxical}.

\cf{Scope and limits stated by the author} Reference points on policy and regulation, legal doctrine and case law are drawn from US experience \pc{nissenbaum2009}{21}{my reference points on policy and regulation, legal doctrine, and case law are drawn from the U.S.}, but the framework is offered as applicable wherever information has context-specific functions and norms; cultural variation is predicted by context-relativity. The cases are those of the time. The framework makes empirical predictions about when protest will occur; future work should check these ``against historical findings'' and develop testable hypotheses \pc{nissenbaum2009}{22}{against historical findings as well as developing testable hypotheses from them}, so the book itself calls the predictive claims untested.

\cf{Limits (card)} Only the introduction and contents are available, so the five parameters (subject, sender, recipient, information type, transmission principle), the norm schema and the heuristic's steps are not shown; the OCR is imperfect; the introduction gives no method for finding norms in practice (surveys, vignettes and policy analysis come from other work); harm is treated as what norms protect against rather than as a typology, so a harm taxonomy has to come from elsewhere (Citron and Solove, Solove); the introduction predates smartphone health apps and uses hospital and nursing-home devices as its health example.

\cf{Use for this thesis} (1) The conceptual basis for describing a privacy-policy statement or app flow as a flow of a specific information type between specific parties under a stated transmission principle, which is what CI-based survey work such as Martin's vignette studies operationalises. For the thesis only the introduction can be cited here; cite Nissenbaum (2004) for the parameters and heuristic. (2) The sentence that norms ``protect people and groups against harm'' links CI to the harm tree: a harm arises when a flow breaks a norm that protected a person in that context. (3) The health-care example (monitoring accepted inside care, resisted in commerce) frames why health apps are a hard case: they sit between a care context and a consumer-technology context, so whose norms apply is itself a coding decision to document. (4) Keep the author's own caveat that the predictive claims are untested when using CI as a lens, not as a measurement.
\end{papercard}
